\section*{\texorpdfstring{\(\mathfrak{sl}_2\) 三元组}{\textbackslash mathfrak\{sl\}\_2 三元组}}

\begin{enumerate}
\def\labelenumi{\arabic{enumi})}
\setcounter{enumi}{32}
\item
  g 中的 \(sl_{2}\) 三元组是指一个元素序列 \((x,h,y)\)——g 中异于 \((0,0,0)\) 且满足 \([h,x]=2x,[h,y]=-2y,[x,y]=-h\) 的。此时 x,y 在 g 中幂零，h 在 g 中半单。
\item
  令 \(x\) 为 \(\mathfrak{g}\) 的非零幂零元。存在 \(h, y \in \mathfrak{g}\) 使 \((x, h, y)\) 是一个 \(\mathfrak{sl}_2\) 三元组。
\item
  令 \((x,h,y)\) 与 \((x',h',y')\) 为两个 \(\mathfrak{sl}_2\) 三元组（在 \(\mathfrak{g}\) 中）。下列条件等价：
\end{enumerate}

\begin{enumerate}
\def\labelenumi{\alph{enumi})}
\item
  存在 \(s \in \text{Aut}_{e}(\mathfrak{g})\) 使 \(sx = x'\)；
\item
  存在 \(s \in \mathrm{Aut}_e(\mathfrak{g})\) 使 \(sx = x', sh = h', sy = y'\)。
\end{enumerate}

\begin{enumerate}
\def\labelenumi{\arabic{enumi})}
\setcounter{enumi}{35}
\tightlist
\item
  若 \(k\) 代数闭，则 35) 的条件 \(a)\) 与 \(b)\) 等价于：\(c)\) 存在 \(s \in \operatorname{Aut}_e(\mathfrak{g})\) 使 \(sh = h'\)。
\end{enumerate}

此外，g 的非零幂零元关于 \(\operatorname{Aut}_{e}(\mathfrak{g})\) 的共轭类个数至多等于 \(3^{l}\)，其中 l 是 g 的秩。

\begin{enumerate}
\def\labelenumi{\arabic{enumi})}
\setcounter{enumi}{36}
\tightlist
\item
  g 的幂零元 x 若其中心化子的维数等于 g 的秩，则称为主幂零元。g 中存在主幂零元。若 k 代数闭，则 g 的主幂零元在 \(\operatorname{Aut}_{e}(\mathfrak{g})\) 下共轭。
\end{enumerate}

